\chapter{Virtual Memory, TLB, Huge Pages, and NUMA}

Every pointer you dereference is a virtual address hardware must translate before any byte is read --- and that translation, the page faults it can trigger, and the non-uniform cost of reaching memory on a multi-socket machine are invisible in source yet decisive for tail latency. This chapter exposes the machinery: page tables and the TLB, demand paging and page faults, huge pages, and NUMA --- and the disciplines that keep it from injecting jitter.

\section*{Chapter Roadmap}
\begin{itemize}
\item 88.1 Why Virtual Memory Exists
\item 88.2 Address Translation, Page Tables, and the TLB
\item 88.3 Page Faults: Minor, Major, and Why They Spike Latency
\item 88.4 Demand Paging and Pre-Faulting
\item 88.5 Huge Pages and TLB Reach
\item 88.6 NUMA: Non-Uniform Memory Access
\item 88.7 The Disciplines for Deterministic Memory
\end{itemize}

\section{Why Virtual Memory Exists}
Virtual memory gives each process a private contiguous address space; the MMU translates virtual to physical on every access. It buys isolation, the illusion of more memory than exists, and simple layout over fragmented physical RAM.

\textbf{Why this matters.} The convenience imposes a tax: every access requires translation, and the mechanisms behind the illusion (demand paging, swapping) can stall the program for milliseconds unpredictably. Latency-critical code cannot opt out but can pin, pre-fault, and use huge pages to make the cost constant and small.

\section{Address Translation, Page Tables, and the TLB}
Memory is divided into 4\,KB pages; the virtual-page-to-frame mapping lives in a multi-level page table (4 levels on x86-64). A page walk costs up to 4 dependent memory accesses. The TLB caches recent translations: a hit is $\sim$1 cycle; a miss triggers a page walk (10s--100s cycles).

\textbf{Why this matters / cost model.} The TLB is small (hundreds to $\sim$1--2K entries); with 4\,KB pages its reach covers only a few MB. Random access over a large footprint thrashes the TLB, paying a page walk per access on top of the data miss --- a hidden tax on hash tables and large graphs that huge pages address.

\section{Page Faults: Minor, Major, and Why They Spike Latency}
A minor fault (page resident, not mapped --- first touch) costs microseconds; a major fault (contents on disk/swap) costs milliseconds with the thread blocked.

\textbf{Why this matters / cost model.} A page fault is an involuntary syscall mid-instruction. A minor fault is why the first write to fresh \texttt{malloc}'d memory is slow (frames are lazily allocated on first touch). A major fault is catastrophic --- one swap-in blows a microsecond budget a thousandfold. Pre-fault and lock memory resident.

\section{Demand Paging and Pre-Faulting}
Default demand paging allocates frames on first access, deferring allocation cost into the hot path.

\begin{lstlisting}[language=C++, caption={Pre-faulting moves minor-fault cost out of the critical path. Linux-specific. Min standard: C++11.}]
void prefault(std::vector<char>& buf) {
    const size_t page = 4096;
    for (size_t i = 0; i < buf.size(); i += page) buf[i] = 0;   // first-touch each page
}
// Or: mmap(..., MAP_POPULATE) pre-faults at mapping time.
\end{lstlisting}

\textbf{Why this matters.} Pre-faulting converts distributed fault latency into one up-front cost during warmup. Combined with warming caches/TLB/predictors (Chapter 106), this makes the first real message as fast as the millionth. The trade-off is startup time and committed memory.

\section{Huge Pages and TLB Reach}
A huge page maps 2\,MB or 1\,GB with one TLB entry instead of 4\,KB, extending TLB reach 512$\times$ (2\,MB).

\begin{lstlisting}[language=bash, caption={Reserving huge pages on Linux (non-portable, privileged).}]
# echo 512 > /proc/sys/vm/nr_hugepages     # 512 * 2MB = 1GB
# mmap(..., MAP_HUGETLB, ...) or Transparent Huge Pages (THP)
\end{lstlisting}

\textbf{Why this matters / cost model.} Huge pages let a multi-GB working set fit in TLB reach, eliminating page walks. Caveats: coarser allocation (fragmentation), higher minor-fault latency (zeroing 2\,MB), and THP's background defrag can introduce jitter --- many latency shops disable THP and use explicit huge pages. A measured optimization, not a default.

\section{NUMA: Non-Uniform Memory Access}
Each socket has local DRAM; local access is fast, remote access crosses an inter-socket link at higher latency and lower bandwidth. Default Linux policy is first-touch: a page lands on the node of the thread that first writes it.

\begin{lstlisting}[language=bash, caption={Binding to a NUMA node; first-touch placement. Non-portable.}]
# numactl --cpunodebind=0 --membind=0 ./app
# A page is placed on the node of the thread that FIRST writes it.
\end{lstlisting}

\textbf{Why this matters / cost model.} Remote access is $\sim$1.5--2$\times$ latency and a fraction of bandwidth. Classic bug: one initialiser thread first-touches a giant array onto one node, then remote workers hammer it. Fix per first-touch: each worker initialises its own memory, and pin threads (Chapter 96). NUMA-obliviousness is a common big-server scalability failure.

\section{The Disciplines for Deterministic Memory}
\begin{itemize}
\item TLB miss --- huge pages, locality, smaller footprint.
\item Minor fault --- pre-fault during warmup.
\item Major fault --- \texttt{mlockall} to lock resident; disable swap.
\item Remote NUMA --- first-touch local; pin threads; \texttt{numactl}.
\item THP defrag --- disable THP; use explicit huge pages.
\end{itemize}

\begin{lstlisting}[language=C++, caption={mlockall pins memory resident, preventing major faults. Linux, privileged. Min standard: C++11.}]
#include <sys/mman.h>
// mlockall(MCL_CURRENT | MCL_FUTURE);
\end{lstlisting}

\textbf{The discipline.} Virtual memory is tuned for average throughput across processes --- wrong for one latency-critical process that cares about the tail. Pre-fault all memory, \texttt{mlock} it, use huge pages, and place memory local to its threads. These turn millisecond jitter into a paid-once cost.
